\chapter{Последовательные интерфейсы: UART, SPI, I$^2$C}
\label{ch:uart}

Ножек у микросхемы ограниченное количество, а данных передавать нужно много. Если
отправлять байт по восьми проводам одновременно (\term{параллельный}
интерфейс), потребуется восемь ножек, восемь дорожек на плате, а на высокой
частоте провода начинают мешать друг другу. Поэтому между микросхемами данные
чаще всего передают \term{последовательно}: по одному проводу, бит за битом. В
этой главе мы разберём три самых распространённых последовательных
интерфейса и построим на ПЛИС один из них --- UART --- целиком, от передатчика
до приёмника.

\section{Три интерфейса}

\begin{table}[H]
\centering
\caption{Сравнение последовательных интерфейсов}
\small
\begin{tabularx}{\textwidth}{l X X X}
\toprule
 & \textbf{UART} & \textbf{SPI} & \textbf{I$^2$C} \\
\midrule
Провода & 2 (TX, RX) + земля & 4 (SCK, MOSI, MISO, CS) & 2 (SDA, SCL) \\
Такт & нет, скорость договорена заранее & есть, ведущий выдаёт SCK & есть, ведущий выдаёт SCL \\
Устройств & 2 (точка---точка) & 1 ведущий и несколько ведомых (по линии CS) & до 127 адресов на 2 проводах \\
Скорость & до единиц Мбит/с & десятки Мбит/с & 0,1--3,4 Мбит/с \\
Типичные применения & консоль, GPS, модули Bluetooth & Flash, дисплеи, АЦП & датчики, часы, EEPROM \\
\bottomrule
\end{tabularx}
\end{table}

\begin{diybox}[Сколько проводов на интерфейс]
\textbf{Задание.} Сравните три интерфейса на готовых примерах и убедитесь, что число сигналов совпадает с табл.~выше.

\textbf{В Logra.} Откройте в карточке «На основе примера» три проекта: «UART-передатчик 8N1», «SPI mode 0: внутренний loopback», «I²C: START и STOP». Вкладка «Код» --- посмотрите список портов модуля верхнего уровня, вкладка «Выводы» --- какие из них нужны платформе.

\textbf{Ожидаемый результат.} У UART отдельные линии передачи и приёма (\code{tx}, \code{rx}) без такта; у SPI есть \code{SCK}, \code{MOSI}, \code{MISO}, \code{CS}; у I$^2$C одна пара линий \code{SDA} и \code{SCL}.

\textbf{Как проверить.} Составьте таблицу «интерфейс --- сигналы --- есть ли такт» и сверьте её с табл.~выше; в SPI-примере обратите внимание, что обмен внутренний (loopback), контакты платы не используются.

\stars{1} повторите: заполните таблицу для трёх проектов. \stars{2} измените: выберите интерфейс для датчика, у которого осталась одна свободная ножка, и обоснуйте. \stars{3} проект: этап «Охранная панель» --- выбор канала связи; критерий приёмки --- в одном абзаце объяснено, почему панель использует UART для связи с компьютером.
\end{diybox}

\section{UART}

\term{UART} (универсальный асинхронный приёмопередатчик) --- самый простой
интерфейс, который стоит в каждом микроконтроллере и в USB-мосте вашей платы
(это именно он связывает ПЛИС с компьютером по USB). Асинхронный --- значит,
\emph{общего тактового сигнала нет}: передатчик и приёмник заранее
договариваются о скорости (бодах), а приёмник сам подстраивается по началу
каждого байта.

В покое линия держится в единице. Байт передаётся кадром: стартовый бит
(0), восемь бит данных \emph{младшим вперёд} и стоповый бит (1). Такой формат
называют 8N1 (8 бит данных, без контроля чётности, 1 стоповый бит).

\begin{figure}[H]
\centering
\begin{tikzpicture}[x=1.0cm, y=0.8cm, font=\sffamily\scriptsize]
  \draw[FpgaBlue, very thick]
    (0,1) -- (1,1) -- (1,0) -- (2,0) -- (2,1) -- (3,1) -- (3,0) -- (8,0) -- (8,1) -- (9,1) -- (9,0) -- (10,0) -- (10,1) -- (11.5,1);
  \foreach \x/\t in {1.5/старт, 2.5/{$d_0$}, 3.5/{$d_1$}, 4.5/{$d_2$}, 5.5/{$d_3$}, 6.5/{$d_4$}, 7.5/{$d_5$}, 8.5/{$d_6$}, 9.5/{$d_7$}, 10.5/стоп}
    \node at (\x,1.7) {\t};
  \foreach \x/\t in {2.5/1, 3.5/0, 4.5/0, 5.5/0, 6.5/0, 7.5/0, 8.5/1, 9.5/0}
    \node[text=FpgaGray] at (\x,-0.5) {\t};
  \foreach \x in {1,2,...,11} \draw[gray!50, dashed] (\x,-0.2) -- (\x,1.4);
  \node[anchor=east] at (0,1) {1};
  \node[anchor=east] at (0,0) {0};
  \node[text=FpgaGray] at (0.5,0.5) {покой};
  \draw[<->] (1,-1.1) -- node[below] {$T_{\text{бит}}$} (2,-1.1);
  \node[font=\sffamily\small] at (6,-2.0) {Кадр для буквы A = \code{0x41} = \code{01000001$_2$}: первым уходит младший бит};
\end{tikzpicture}
\caption{Кадр UART 8N1}
\label{fig:uart_frame}
\end{figure}

\subsection{Скорость передачи}

Время одного бита $T_{\text{бит}}$ равно $1/\text{скорость}$. Для
115\,200 бод оно составляет около $8{,}7$~мкс. ПЛИС работает на 27~МГц, поэтому
на один бит приходится
\[
\text{CLKS\_PER\_BIT} = \frac{27\,000\,000}{115\,200} \approx 234 \text{ такта}.
\]
И передатчик, и приёмник считают такты и меняют (или читают) бит каждые
\code{CLKS\_PER\_BIT} тактов.

\subsection{Передатчик}

Передатчик проще: он сохраняет собранный кадр в сдвиговом регистре и по
таймеру выдаёт его биты на линию. Между байтами линия остаётся в покое (1).
Автомат состоит из двух состояний: \code{IDLE}, где он ждёт сигнал
\code{start}, и \code{SEND}, где передаются десять бит кадра (старт, 8 данных, стоп).

\subsection{Приёмник}

Приёмник сложнее, потому что он не знает, когда начнётся кадр, и общего
такта у него нет. Он делает так:

\begin{enumerate}
  \item Ждёт, пока линия упадёт в 0: это начало стартового бита.
  \item Отсчитывает \emph{полбита}, чтобы оказаться в середине стартового бита, и
        проверяет, что линия всё ещё в нуле (иначе это была помеха).
  \item Дальше читает линию строго через один бит --- всегда в середине бита, где
        сигнал устоялся дальше всего от краёв.
  \item Читает стоповый бит: он обязан быть единицей, иначе кадр считается
        повреждённым.
\end{enumerate}

Кроме того, вход \code{rx} приходит снаружи и никак не связан с тактом ПЛИС,
поэтому на входе приёмника стоят два последовательных триггера
(синхронизатор, см. главу~\ref{ch:clocks}).

\vfile{\codepath/08_uart/uart.v}{code/08\_uart/uart.v}

Файл содержит три модуля. \code{uart\_tx} и \code{uart\_rx} --- сами
приёмопередатчики с параметром \code{CLKS\_PER\_BIT}, а \code{uart\_demo} --- их
демонстрация: передатчик раз за разом отправляет буквы от A до Z, а выход
\code{tx} подан на вход приёмника этого же модуля, поэтому приёмник
получает то, что отправлено. Для просмотра в симуляторе скорость
сделана очень высокой (\code{CLKS\_PER\_BIT = 8}), иначе пришлось бы ждать
тысячи тактов.

\shot{n-uart-wave}{Передача буквы A в Logra: сначала на \code{tx} стартовый бит и восемь бит данных, затем приёмник выставляет \code{rx\_valid} и в \code{rx\_data} появляется \code{0x41}}{fig:n-uart-wave}

На диаграмме (рис.~\ref{fig:n-uart-wave}) можно пересчитать кадр: линия опускается
(старт), несколько бит отправляются, и в конце, ровно через 10 периодов бита,
приёмник сообщает о готовом байте импульсом \code{rx\_valid} и выставляет
\code{rx\_data}.

\begin{infobox}[Как подключить к компьютеру]
USB-мост на вашей плате выводит UART на компьютер. Достаточно задать
\code{CLKS\_PER\_BIT = 234} (115\,200 бод при 27~МГц), назначить ножки \code{tx} и
\code{rx} в планировщике выводов, загрузить схему и открыть вкладку «UART-терминал»
в разделе «Подключение»: там будут появляться отправленные буквы.
\end{infobox}

\begin{diybox}[Передать букву A]
\textbf{Задание.} Передайте по UART букву A (\code{0x41}), разберите кадр по диаграмме и пропустите его через приёмник.

\textbf{В Logra.} Откройте «UART-передатчик 8N1» и «UART-приёмник 8N1» либо \code{uart.v} из книги (\code{uart\_demo}, \code{CLKS\_PER\_BIT = 8}). Вкладка «Диаграммы»: сигналы \code{tx}, \code{rx\_valid}, \code{rx\_data}.

\textbf{Ожидаемый результат.} Линия опускается на старт, затем идут биты $1,0,0,0,0,0,1,0$ (младший первым), затем стоповая единица; приёмник поднимает \code{rx\_valid} и выставляет \code{rx\_data = 0x41}.

\textbf{Как проверить.} Диаграммы (рис.~\ref{fig:n-uart-wave}): пересчитайте биты кадра, убедитесь, что ширина бита равна \code{CLKS\_PER\_BIT} тактов, а \code{rx\_valid} приходит в конце кадра. Для платы: задайте \code{CLKS\_PER\_BIT = 234}, назначьте выводы на вкладке «Выводы», вкладка «Подключение» $\to$ «UART-терминал».

\stars{1} повторите: буква A на диаграмме. \stars{2} измените: передайте \code{0x55} и убедитесь, что линия строго чередуется. \stars{3} проект: этап «Охранная панель» --- отправка состояния по UART; критерий приёмки --- при каждой смене \code{state} в терминале появляется байт (0, 1 или 2) на 115\,200 бод.
\end{diybox}

\section{SPI}

В \term{SPI} есть общий такт \code{SCK}, и поэтому он проще и быстрее. Ведущий
выдаёт тактовый сигнал и данные на линию \code{MOSI} (master out, slave in) и
одновременно получает данные от ведомого по линии \code{MISO}. Выбирает ведомого
сигнал \code{CS} (chip select) с активным низким уровнем.

\begin{figure}[H]
\centering
\begin{tikzpicture}[x=0.7cm, y=0.8cm, font=\sffamily\scriptsize]
  \node[anchor=east] at (-0.2,2.7) {CS};
  \draw[FpgaBlue, very thick] (0,3) -- (1,3) -- (1,2.4) -- (17,2.4) -- (17,3) -- (18,3);
  \node[anchor=east] at (-0.2,1.5) {SCK};
  \draw[FpgaBlue, very thick] (0,1.2) -- (2,1.2);
  \foreach \i in {0,...,7}
    \draw[FpgaBlue, very thick] (2+\i*2,1.2) -- (2+\i*2,1.8) -- (3+\i*2,1.8) -- (3+\i*2,1.2) -- (4+\i*2,1.2);
  \node[anchor=east] at (-0.2,0) {MOSI};
  \draw[FpgaGreen!60!black, very thick] (0,-0.3) -- (1,-0.3) -- (1,0.3) -- (3,0.3) -- (3,-0.3) -- (5,-0.3) -- (5,0.3) -- (9,0.3) -- (9,-0.3) -- (13,-0.3) -- (13,0.3) -- (15,0.3) -- (15,-0.3) -- (17,-0.3);
  \foreach \i/\v in {0/1, 1/0, 2/1, 3/1, 4/0, 5/0, 6/1, 7/0} \node at (2+\i*2,-1.0) {\v};
  \foreach \i in {0,...,7} \draw[gray!50, dashed] (2+\i*2,-0.5) -- (2+\i*2,1.9);
  \node[align=left, anchor=west] at (0,-2.1) {Данные меняются при спаде SCK и читаются ведомым по нарастающему фронту (пунктир), режим 0};
\end{tikzpicture}
\caption{Передача байта \code{10110010} по SPI. Данные меняются на спаде такта и читаются на нарастающем фронте}
\end{figure}

У SPI четыре \term{режима} --- сочетания полярности тактового сигнала в
покое (CPOL) и фазы (CPHA, читают по первому или второму фронту). Режимы
нужно сверять с документацией ведомого устройства. Простейший передатчик SPI
в ПЛИС --- это сдвиговый регистр из главы~\ref{ch:regs}, который сдвигается на
каждом такте SCK, и делитель частоты, формирующий сам SCK.

\begin{diybox}[SPI: смена и выборка]
\textbf{Задание.} Отправьте байт \code{0xA6} по SPI в режиме 0 и проверьте, что ведущий получил то же самое по внутреннему loopback.

\textbf{В Logra.} Проект «SPI mode 0: внутренний loopback» (карточка «На основе примера»), параметр \code{HALF\_TICKS = 1}. Вкладка «Диаграммы»: \code{sck}, \code{mosi}, \code{cs\_n}, \code{received}.

\textbf{Ожидаемый результат.} \code{SCK} в покое равен 0; данные на \code{MOSI} меняются на спаде, читаются на нарастающем фронте; после восьми тактов \code{received = 0xA6}.

\textbf{Как проверить.} Диаграммы: сравните порядок битов \code{MOSI} с \code{0xA6 = 10100110$_2$} (старший первым) и совпадение с рис. кадра SPI.

\stars{1} повторите: \code{0xA6}. \stars{2} измените: отправьте \code{0xC3} и предскажите значение \code{received}. \stars{3} проект: этап «Охранная панель» --- байт состояния по SPI; критерий приёмки --- \code{cs\_n} опущен ровно на 8 тактов \code{SCK}, данные не меняются на нарастающем фронте.
\end{diybox}

\section{I$^2$C}

\term{I$^2$C} использует всего два провода, но взамен усложняет протокол.
\code{SDA} несёт данные, \code{SCL} --- такт. Обе линии выполнены с открытым
стоком и подтянуты резисторами к питанию (глава~\ref{ch:electrical}), поэтому
любое устройство может опустить линию в 0, а единицу создаёт подтяжка.

Обмен начинается с условия \term{START} (SDA падает при высоком SCL), потом
идёт 7-битный адрес устройства, бит направления (чтение или запись) и
подтверждение \term{ACK}: адресованная микросхема сама опускает SDA на
девятом такте. Затем передаются байты, каждый с подтверждением, и в конце
условие \term{STOP} (SDA поднимается при высоком SCL). На одной линии может сидеть
много устройств, у каждого свой адрес, а столкновений не бывает именно
благодаря открытому стоку.

Проект с I$^2$C сложнее: нужен автомат из десятков состояний, а настоящее
устройство подтверждает приём. В лаборатории Logra есть модели OLED-дисплея,
часов и памяти на I$^2$C, которые отвечают побитно, так что можно проверить
свой автомат до включения платы.

\begin{diybox}[I\texorpdfstring{$^2$}{2}C: START и STOP]
\textbf{Задание.} Найдите на диаграмме условия START и STOP и убедитесь, что линия с открытым стоком не может быть принудительно подтянута к единице.

\textbf{В Logra.} Проект «I²C: START и STOP» (карточка «На основе примера»), \code{STEP\_TICKS = 1}; вкладка «Диаграммы»: \code{sda}, \code{scl}, \code{drive\_sda\_low}.

\textbf{Ожидаемый результат.} START: \code{SDA} падает при высоком \code{SCL}; затем один учебный бит; STOP: \code{SDA} растёт при высоком \code{SCL}. Адреса и ACK в этом примере нет.

\textbf{Как проверить.} Диаграммы: отметьте оба условия. Затем откройте вкладку «Лаборатория» (группа «Проверить»), в списке деталей найдите I$^2$C-устройство и посмотрите, как оно отвечает на адрес.

\stars{1} повторите: найдите START и STOP. \stars{2} измените: добавьте в последовательность 7 бит адреса и бит ACK. \stars{3} проект: этап «Охранная панель» --- индикатор на I$^2$C; критерий приёмки --- на диаграмме видны START, адрес, бит направления, ACK, байт данных, STOP.
\end{diybox}

\section{Вопросы и задачи}
\label{sec:tasks-uart}

Ответы собраны в приложении~\ref{app:answers}.

\begin{taskbox}[Задачи к главе]
\begin{enumerate}
  \item \stars{1}~Какой \code{CLKS\_PER\_BIT} нужен для 9600 бод при частоте 27~МГц?
  \item \stars{1}~Сколько времени займёт передача одного байта в формате 8N1 на скорости
        115\,200 бод? Сколько байт в секунду можно передать?
  \item \stars{1}~Нарисуйте кадр UART для байта \code{0x55} (\code{01010101}).
  \item \stars{2}~Почему приёмник UART читает биты в \emph{середине} их длительности, а не в начале?
  \item \stars{2}~Тактовая частота приёмника на 3\,\% выше, чем у передатчика. Не сломается
        ли приём кадра из 10 бит? Какое расхождение скоростей допустимо?
  \item \stars{2}~Чем SPI быстрее UART и в чём его недостаток?
  \item \stars{2}~Найдите ошибку. Приёмник UART после опускания линии сразу читает биты, отсчитывая
        \code{CLKS\_PER\_BIT} тактов от первого фронта. На скорости 115\,200 бод с хорошим
        кабелем всё работает, а на длинном проводе появляются сбои. Почему? Как исправить?
  \item \stars{3}~Проект. Модуль передаёт каждую секунду байт состояния панели (0, 1 или 2)
        на скорости 115\,200 бод. Критерий приёмки: на диаграмме кадр состоит из
        старта, 8 бит младшим вперёд и стопа, а в терминале раз в секунду приходит символ.
  \item \stars{3}~Проект. Свяжите два модуля \code{uart\_tx} и \code{uart\_rx} в одном проекте
        и покажите принятый байт на светодиодах. Критерий приёмки: все 256 значений проходят
        без искажений, подача ошибочного стопового бита не даёт \code{rx\_valid}.
\end{enumerate}
\end{taskbox}
